\documentclass{article}

% Submission mode (no option) anonymises the authors and turns on the
% line numbers down the margin, which is the most distinctive thing this
% template renders. `preprint` or `final` would switch both off.
\usepackage{neurips_2026}

\usepackage[utf8]{inputenc}
\usepackage[T1]{fontenc}
\usepackage{hyperref}
\usepackage{url}
\usepackage{booktabs}
\usepackage{amsfonts}
\usepackage{nicefrac}
\usepackage{microtype}

\title{Ink Density Saturation in Duan Inkstones}

\author{%
  David S.~Hippocampus\thanks{Placeholder author from the NeurIPS template;
    submission mode hides this block.} \\
  Department of Inkstone Studies\\
  Cranberry-Lemon University\\
  Pittsburgh, PA 15213 \\
  \texttt{hippo@cs.cranberry-lemon.edu} \\
}

\begin{document}

\maketitle

\begin{abstract}
  A minimal document in the NeurIPS 2026 format. The content mirrors
  \texttt{latex-rendering-check.tex}: ink density saturates as a function
  of grind time, and the stone --- not the stick --- sets the ceiling.
\end{abstract}

\section{Model}

Let $D(t)$ denote ink density after $t$ minutes of grinding. Density
follows a saturating exponential,
\begin{equation}
  \label{eq:density}
  D(t) = D_{\max}\left(1 - e^{-t/\tau}\right),
\end{equation}
where $\tau$ is the stone's time constant and $D_{\max}$ its ceiling. The
relationship was established empirically by \citet{tanaka2019}, and half
saturation is reached at $\nicefrac{\ln 2}{1} \cdot \tau$.

\section{Results}

Table~\ref{tab:stones} lists time constants for three stones. The
grinding surface, not the ink stick, dominates \citep{mizuta1968}.

\begin{table}[t]
  \caption{Saturation parameters by stone.}
  \label{tab:stones}
  \centering
  \begin{tabular}{lrr}
    \toprule
    Stone & $\tau$ (min) & $D_{\max}$ \\
    \midrule
    Duan (Laokeng)   & 4.1 & 0.94 \\
    Duan (Majiakeng) & 6.8 & 0.87 \\
    She              & 9.2 & 0.81 \\
    \bottomrule
  \end{tabular}
\end{table}

\section{Checklist}

Opening this file and pressing preview should produce a single-column PDF
with line numbers down the left margin, the NeurIPS footer on page one, a
numbered equation, a booktabs table, and resolved citations to
\citet{tanaka2019} and \citet{chen2021}.

\bibliographystyle{plainnat}
\bibliography{references}

\end{document}
